% parte2-coro-master.tex -- choir placeholder and master.
\chapter{Coro e master}\label{cap:coro-master}
Il coro (\texttt{pitchDetectorChoirMcAdams}) suona in quattro istanze, una per canale.
Ogni istanza ha 16 bande: l'inviluppo della banda $k$ intorno a $f \cdot k$ modula un rumore filtrato intorno a $f \cdot k^{a}$, con $f$ = 48, 48, 96 e \qty{96}{\hertz}, $a$ = 1; 1,01; 1,1; 0,9, fattore di merito 350 e rilascio di \qty{1.5}{\second}.
Lo studio del coro seguirà quello di stunedrev: la scheda qui sotto tiene il suo posto nel sistema.

\scheda{coro-uscita}{Uscita del coro (CC86)}
{}
{}
{non ancora studiato: lo studio seguirà quello di stunedrev}
{}
{\daprovare}
{}

\section{Master}
Il master chiude la catena e porta i quattro canali alle uscite 9--12, verso STONED.

\scheda{master}{Master (CC88)}
{fader lineare, levigato da \texttt{si.smoo} della libreria standard}
{come nel patch}
{nessuna oltre al livello}
{invariato; nel porting gli stadi di guadagno del patch diventano i fader delle tracce di Reaper}
{\daprovare}
{log di sessione, blocco 6 di delRM}
